% ============================================================
% BAB 2 PENDAHULUAN
% ============================================================
\chapter{PENDAHULUAN}

\section{Latar Belakang}
% → masalah penelitian muncul

\section{Rumusan Masalah}
% → apa yang harus dijawab

\section{Tujuan Penelitian}
% → apa yang akan dilakukan untuk menjawab rumusan masalah

\section{Batasan Penelitian}
% → ruang lingkup penelitian

\section{Kontribusi Penelitian}
% Apa yang penelitian ini hasilkan?
% bisa berupa :
% benchmark dataset semantic segmentation NLB;
% model lightweight semantic segmentation;
% framework cross-dataset evaluation;
% metode disease severity estimation;

\section{Originalitas Penelitian}
% Apa yang belum pernah dilakukan penelitian sebelumnya dan dilakukan dalam penelitian ini?

\section{Manfaat Penelitian}
% → dampak teoretis dan praktis
\subsection{Manfaat Teoritis}
\subsection{Manfaat Praktis}

\section{Kerangka dan Peta Jalan Penelitian}
\subsection{Fishbone Diagram Penelitian}
% Fishbone menunjukkan hubungan:
% Research Problem → Limitations → Research Gap → RQ → Objectives → Contributions

\subsection{Peta Jalan Penelitian}
% sedangkan roadmap menunjukkan perjalanan penelitian:
% Conference 1 → Conference 2 → Journal 1 → ... → Journal 4 → Disertasi

\section{Sistematika Penulisan}



Penelitian ini bertujuan untuk mengembangkan suatu pendekatan yang dapat meningkatkan efektivitas dan efisiensi proses analisis data. Perkembangan teknologi yang semakin pesat telah mendorong penggunaan berbagai metode komputasi untuk menyelesaikan permasalahan yang kompleks.

\begin{figure}[htbp]
    \centering
    \includegraphics[width=0.5\textwidth]{1_kotak_orange}
    \caption{Laju infiltrasi permukaan tanah alluvial yang diobservasi
    setiap jarak interval 25-ft di Rio Grande dekat Socorro, New Mexico
    (Gelhar, 1993; data diperoleh dari Gelhar et al., 1983)}
   \label{fig:laju-infiltrasi}
\end{figure}
\begin{figure}[htbp]
    \centering
    \includegraphics[width=0.5\textwidth]{1_kotak_orange}
    \caption{Laju infiltrasi permukaan tanah alluvial yang diobservasi
    setiap jarak interval 25-ft di Rio Grande dekat Socorro, New Mexico
    (Gelhar, 1993; data diperoleh dari Gelhar et al., 1983)}
   \label{fig:laju-infiltrasi1}
\end{figure}
Penelitian ini bertujuan untuk mengembangkan suatu pendekatan yang dapat meningkatkan efektivitas dan efisiensi proses analisis data. Perkembangan teknologi yang semakin pesat telah mendorong penggunaan berbagai metode komputasi untuk menyelesaikan permasalahan yang kompleks.

Penelitian ini bertujuan untuk mengembangkan suatu pendekatan yang dapat meningkatkan efektivitas dan efisiensi proses analisis data. Perkembangan teknologi yang semakin pesat telah mendorong penggunaan berbagai metode komputasi untuk menyelesaikan permasalahan yang kompleks.


\section{Rumusan Masalah}
Berdasarkan latar belakang tersebut, permasalahan utama dalam penelitian
ini adalah belum tersedianya sistem...
Permasalahan tersebut dirumuskan ke dalam beberapa pertanyaan penelitian sebagai berikut :

\begin{enumerate}[
    label=\arabic*.,
        leftmargin=1cm, %leftmargin=*, % ini jika margin = 0
    itemsep=0.3\baselineskip, % jarak antar item
    topsep=0pt,
    parsep=0pt
]
\item Bagaimana merancang dataset dataset dataset dataset dataset dataset dataset dataset dataset dataset dataset dataset dataset dataset dataset dataset dataset dataset dataset dataset dataset dataset dataset dataset?

\item Bagaimana mengembangkan model model model model model model model model model model model model model model model model model model model model model model model model model model model model model model model model model?

\item Bagaimana mengimplementasikan sistem sistem sistem sistem sistem sistem sistem sistem sistem sistem sistem sistem sistem sistem sistem sistem sistem sistem sistem sistem sistem sistem ?

\item Bagaimana mengevaluasi performa sistem dari aspekA aspekA aspekA aspekA aspekA aspekA aspekA aspekA aspekA aspekA aspekA aspekA aspekA aspekA aspekA aspekA aspekA aspekA aspekA aspekA aspekA aspekA aspekA aspekA aspekA?
    
\end{enumerate}

\begin{figure}[htbp]
    \centering
    \includegraphics[width=0.5\textwidth]{1_kotak_orange}
    \caption{Laju infiltrasi permukaan tanah alluvial yang diobservasi
    setiap jarak interval 25-ft di Rio Grande dekat Socorro, New Mexico
    (Gelhar, 1993; data diperoleh dari Gelhar et al., 1983)}
   \label{fig:laju-infiltrasi2}
\end{figure}

Penelitian ini bertujuan untuk mengembangkan suatu pendekatan yang dapat meningkatkan efektivitas dan efisiensi proses analisis data. Perkembangan teknologi yang semakin pesat telah mendorong penggunaan berbagai metode komputasi untuk menyelesaikan permasalahan yang kompleks.

\begin{table}[H]
    \centering
    \caption{Nilai Tipikal dari Batas Cair, Batas Plastis, dan Aktivitas 
        dari Beberapa Mineral Lempung}
    \label{tab:nilai-tipikal-mineral-lempung}
    \renewcommand{\arraystretch}{1.15} % Ini mengatur tinggi baris tabel. Nilai default kira-kira 1.0.
    \setlength{\tabcolsep}{6pt} % mengatur jarak horizontal antara isi sel dengan garis/batas kolom.

    \begin{tabular}{
        |>{\raggedright\arraybackslash}m{4.3cm}
        |>{\centering\arraybackslash}m{2.0cm}
        |>{\centering\arraybackslash}m{2.0cm}
        |>{\centering\arraybackslash}m{2.7cm}|
    }
        \hline
        
        \multicolumn{1}{|>{\centering\arraybackslash}m{4.3cm}|}
        {\textbf{Mineral}} &
        \textbf{Batas Cair, LL} &
        \textbf{Batas Plastis, PL} &
        \textbf{Aktivitas, A} \\
        \hline

        \textit{Kaolinite} &
        35--100 &
        20--40 &
        0,3--0,5 \\
        \hline

        \textit{Illite} &
        60--120 &
        35--60 &
        0,5--1,2 \\
        \hline

        \textit{Montmorillonite} &
        100--900 &
        50--100 &
        1,5--7,0 \\
        \hline

        \textit{Halloysite (hydrated)} &
        50--70 &
        40--60 &
        0,1--0,2 \\
        \hline

        \textit{Halloysite (dehydrated)} &
        40--55 &
        30--45 &
        0,4--0,6 \\
        \hline

        \textit{Attapulgite} &
        150--250 &
        100--125 &
        0,4--1,3 \\
        \hline

        \textit{Allophane} &
        200--250 &
        120--150 &
        0,4--1,3 \\
        \hline
    \end{tabular}

    \vspace{0.5\baselineskip}

    \begin{minipage}{1\textwidth}
        \raggedright
        Sumber: Das, 2002
    \end{minipage}

\end{table}

Penelitian ini bertujuan untuk mengembangkan suatu pendekatan yang dapat meningkatkan efektivitas dan efisiensi proses analisis data. Perkembangan teknologi yang semakin pesat telah mendorong penggunaan berbagai metode komputasi untuk menyelesaikan permasalahan yang kompleks.

\section{Tujuan Penelitian}
Penelitian ini bertujuan untuk mengembangkan suatu pendekatan yang dapat meningkatkan efektivitas dan efisiensi proses analisis data. Perkembangan teknologi yang semakin pesat telah mendorong penggunaan berbagai metode komputasi untuk menyelesaikan permasalahan yang kompleks.

Penelitian ini bertujuan untuk mengembangkan suatu pendekatan yang dapat meningkatkan efektivitas dan efisiensi proses analisis data. Perkembangan teknologi yang semakin pesat telah mendorong penggunaan berbagai metode komputasi untuk menyelesaikan permasalahan yang kompleks.

\subsection{Tujuan Penelitian}
Penelitian ini bertujuan untuk mengembangkan suatu pendekatan yang dapat meningkatkan efektivitas dan efisiensi proses analisis data. Perkembangan teknologi yang semakin pesat telah mendorong penggunaan berbagai metode komputasi untuk menyelesaikan permasalahan yang kompleks.

Penelitian ini bertujuan untuk mengembangkan suatu pendekatan yang dapat meningkatkan efektivitas dan efisiensi proses analisis data. Perkembangan teknologi yang semakin pesat telah mendorong penggunaan berbagai metode komputasi untuk menyelesaikan permasalahan yang kompleks.

\subsection{Tujuan Penelitian}
Penelitian ini bertujuan untuk mengembangkan suatu pendekatan yang dapat meningkatkan efektivitas dan efisiensi proses analisis data. Perkembangan teknologi yang semakin pesat telah mendorong penggunaan berbagai metode komputasi untuk menyelesaikan permasalahan yang kompleks.

Penelitian ini bertujuan untuk mengembangkan suatu pendekatan yang dapat meningkatkan efektivitas dan efisiensi proses analisis data. Perkembangan teknologi yang semakin pesat telah mendorong penggunaan berbagai metode komputasi untuk menyelesaikan permasalahan yang kompleks.

\section{Batasan Penelitian}
Penelitian ini bertujuan untuk mengembangkan suatu pendekatan yang dapat meningkatkan efektivitas dan efisiensi proses analisis data. Perkembangan teknologi yang semakin pesat telah mendorong penggunaan berbagai metode komputasi untuk menyelesaikan permasalahan yang kompleks.

Penelitian ini bertujuan untuk mengembangkan suatu pendekatan yang dapat meningkatkan efektivitas dan efisiensi proses analisis data. Perkembangan teknologi yang semakin pesat telah mendorong penggunaan berbagai metode komputasi untuk menyelesaikan permasalahan yang kompleks.

\section{Kontribusi Penelitian}
Penelitian ini bertujuan untuk mengembangkan suatu pendekatan yang dapat meningkatkan efektivitas dan efisiensi proses analisis data. Perkembangan teknologi yang semakin pesat telah mendorong penggunaan berbagai metode komputasi untuk menyelesaikan permasalahan yang kompleks.

Penelitian ini bertujuan untuk mengembangkan suatu pendekatan yang dapat meningkatkan efektivitas dan efisiensi proses analisis data. Perkembangan teknologi yang semakin pesat telah mendorong penggunaan berbagai metode komputasi untuk menyelesaikan permasalahan yang kompleks.

Penelitian sebelumnya telah membahas permasalahan tersebut
\parencite{neuman1980}.

Pendekatan lain juga telah dikembangkan
\parencite{siregar2006}.

\section{Sistematika Penulisan}
Penelitian ini bertujuan untuk mengembangkan suatu pendekatan yang dapat meningkatkan efektivitas dan efisiensi proses analisis data. Perkembangan teknologi yang semakin pesat telah mendorong penggunaan berbagai metode komputasi untuk menyelesaikan permasalahan yang kompleks.

Penelitian ini bertujuan untuk mengembangkan suatu pendekatan yang dapat meningkatkan efektivitas dan efisiensi proses analisis data. Perkembangan teknologi yang semakin pesat telah mendorong penggunaan berbagai metode komputasi untuk menyelesaikan permasalahan yang kompleks.

\section{Sistematika Penulisan}
Penelitian ini bertujuan untuk mengembangkan suatu pendekatan yang dapat meningkatkan efektivitas dan efisiensi proses analisis data. Perkembangan teknologi yang semakin pesat telah mendorong penggunaan berbagai metode komputasi untuk menyelesaikan permasalahan yang kompleks.

\section{Sistematika Penulisan}
Penelitian ini bertujuan untuk mengembangkan suatu pendekatan yang dapat meningkatkan efektivitas dan efisiensi proses analisis data. Perkembangan teknologi yang semakin pesat telah mendorong penggunaan berbagai metode komputasi untuk menyelesaikan permasalahan yang kompleks.

\section{Sistematika Penulisan}
Penelitian ini bertujuan untuk mengembangkan suatu pendekatan yang dapat meningkatkan efektivitas dan efisiensi proses analisis data. Perkembangan teknologi yang semakin pesat telah mendorong penggunaan berbagai metode komputasi untuk menyelesaikan permasalahan yang kompleks.

\section{Sistematika Penulisan}
Penelitian ini bertujuan untuk mengembangkan suatu pendekatan yang dapat meningkatkan efektivitas dan efisiensi proses analisis data. Perkembangan teknologi yang semakin pesat telah mendorong penggunaan berbagai metode komputasi untuk menyelesaikan permasalahan yang kompleks.

\section{Sistematika Penulisan}
Penelitian ini bertujuan untuk mengembangkan suatu pendekatan yang dapat meningkatkan efektivitas dan efisiensi proses analisis data. Perkembangan teknologi yang semakin pesat telah mendorong penggunaan berbagai metode komputasi untuk menyelesaikan permasalahan yang kompleks.

Penelitian sebelumnya telah membahas permasalahan tersebut
\parencite{neuman1980}.

Pendekatan lain juga telah dikembangkan
\parencite{siregar2006}.

Menurut \parencite{santoso2024deep}, deep learning dapat digunakan
untuk mendeteksi penyakit tanaman.

Metode segmentasi juga dikembangkan oleh
\parencite{santoso2024semantic}.

Sementara itu, pendekatan lightweight dibahas oleh
\parencite{santoso2024lightweight}.

Penelitian sebelumnya dapat dilihat pada
\parencite{santoso2023unet}.

Pendekatan lain dikembangkan oleh
\textcite{wijaya2024cnn}.

aaaa \parencite{todd2005groundwater}

proseding \parencite{neuman1980adjoint}

